\subsection{Stage 1 --- Configuration Generation}
\label{sec:stage1}

The configuration generator produces randomized alloy compositions over the
twelve elements covered by at least one of the MEAM library files currently
available to the project: Al, Co, Cr, Cu, Fe, Mg, Mn, Mo, Ni, Si, Ti and Zn. Each generated configuration is a JSON object containing the
species composition (mass fractions summing to exactly unity at $10^{-4}$
precision), the simulation box edge in metres, the target temperature and
the requested number of MD steps. A representative example, corresponding
to a simplified Al~7075 composition, is shown below:

\begin{verbatim}
{
  "composition": { "Al": 0.91, "Mg": 0.029, "Zn": 0.061 },
  "box_size_m": 4e-09,
  "temperature": 0.1,
  "total_steps": 1000
}
\end{verbatim}

To bias the dataset towards practically interesting alloys, aluminium is
forced to be the dominant species whenever it is present in the random
draw. Configurations corresponding to commercial alloys (Al~2024, Al~2219,
Al~5052, Al~7075) are included as named reference entries in addition to
the random ones. At the time of this report, $7924$ configuration
records have been generated, providing a wide population from which
Stage~2 selects samples for MD simulation. The element-occurrence
counts across this set are summarised in
Figure~\ref{fig:element_freq}: aluminium appears in $\sim 45\,\%$ of
configurations (consistent with the Al-dominance bias), while every
other element is sampled in roughly $30\!-\!37\,\%$ of records, giving
the merge logic of Stage~4 enough cross-pair coverage to seed every
binary interaction.

\begin{figure}[h]
	\centering
	\includegraphics[width=0.85\textwidth]{figures/element_frequency.png}
	\caption{Number of configurations in which each element is present,
	out of the $7924$ records currently in the dataset. The dominance of
	Al reflects the bias rule that retains aluminium as the majority
	species whenever it is drawn.}
	\label{fig:element_freq}
\end{figure}
